\section*{Abstract}
		The use of natural units in theoretical physics is a fundamental concept that can be comprehensively explained and contextualized within the framework of T0 theory. This treatise illuminates the principle of dimensional reduction, the advantages for calculations, the particular relevance for T0 theory, and the necessity of explicit SI units in practice. Finally, it emphasizes the deeper insight that physics ultimately rests on dimensionless geometric relationships.

	
	
	\section{Basic Principle of Natural Units}
	\label{sec:grundprinzip}
	
	\subsection{The Principle of Dimensional Reduction}
	In natural units, one sets fundamental constants to 1:
	\begin{itemize}
		\item \textbf{Speed of light}: $c = 1$
		\item \textbf{Reduced Planck constant}: $\hbar = 1$
		\item \textbf{Boltzmann constant}: $k_B = 1$
		\item \textbf{Sometimes}: $G = 1$ (Planck units)
	\end{itemize}
	
	\subsection{Mathematical Consequence}
	This does not mean that these constants ``disappear,'' but that they serve as \textbf{scale setters}:
	\begin{equation}
		E = m c^2 \quad \Rightarrow \quad E = m \quad \text{(since $c=1$)}
	\end{equation}
	\begin{equation}
		E = \hbar \omega \quad \Rightarrow \quad E = \omega \quad \text{(since $\hbar=1$)}
	\end{equation}
	
	\section{Advantages for Calculations}
	
	\subsection{Simplified Formulas}
	\textbf{With SI units:}
	\begin{equation}
		E = \sqrt{(p c)^2 + (m c^2)^2}
	\end{equation}
	\textbf{In natural units:}
	\begin{equation}
		E = \sqrt{p^2 + m^2}
	\end{equation}
	
	\subsection{Transparent Dimensional Analysis}
	All quantities can be traced back to one fundamental dimension (typically energy):
	\begin{table}[h]
		\centering
		\adjustbox{max width=\textwidth}{%
\begin{tabular}{lll}
			\toprule
			\textbf{Quantity} & \textbf{Natural Dimension} & \textbf{SI Equivalent} \\
			\midrule
			Length & $[E]^{-1}$ & $\hbar c / E$ \\
			Time & $[E]^{-1}$ & $\hbar / E$ \\
			Mass & $[E]$ & $E/c^2$ \\
			\bottomrule
		\end{tabular}%
}
		\caption{Dimensional relationships in natural units}
	\end{table}
	
	\section{Particular Relevance in T0 Theory}
	
	\subsection{Geometric Nature of Constants}
	T0 theory shows particularly clearly why natural units are fundamental:
	\begin{equation}
		\alpha = \xi \cdot \left( \frac{E_0}{1~\mathrm{MeV}} \right)^2
	\end{equation}
	This makes explicit that the fine structure constant is a \textbf{purely dimensionless geometric relationship}.
	
	\subsection{The $\xi$-Parameter as Fundamental Geometry Factor}
	The derivation:
	\begin{equation}
		\xi = \frac{4}{3} \times 10^{-4}
	\end{equation}
	is intrinsically dimensionless and represents the fundamental space geometry -- independent of human units of measurement.
	
	\textbf{Important:} $\xi$ alone is not directly equal to $1/m_e$ or $1/E$, but requires specific scaling factors for different physical quantities.
	
	\section{The Scaling Factor $S_{T0}$}
	\label{sec:scaling-derivation}
	
	The scaling factor connects the electron mass in T0 units with its SI value:
	
	\begin{align}
		\textbf{1. T0 units:} \quad & m_e^{\mathrm{T0}} = 0.511 \quad \text{(electron mass in MeV)} \\
		\textbf{2. Measurement:} \quad & m_e^{\mathrm{SI}} = 9.109 \times 10^{-31}~\mathrm{kg} \\
		\textbf{3. Scaling:} \quad & S_{T0} = m_e^{\mathrm{SI}}/m_e^{\mathrm{T0}} = 1.7827\!\times\!10^{-30} \\
		\textbf{4. Comparison:} \quad & 1~\mathrm{MeV}/c^2 = 1.7827 \times 10^{-30}~\mathrm{kg}
	\end{align}
	
	The number $0.511$ is the electron mass in MeV, i.e.\ the measured value in a human-chosen unit; no $\xi$ calculation giving $0.511$ appears in this document. The corpus ladder $m_e = \frac43\,\xi^{3/2}\,v$ needs $v$ as a scale and gives $0.505$ with $v$ in MeV. Hence $S_{T0} = m_e^{\mathrm{SI}}/0.511 = 1~\mathrm{MeV}/c^2$ by construction and not a prediction of the MeV scale: $S_{T0}$ is the conversion between T0 units (masses in MeV) and SI units.
	
	\section{Quantized Mass Calculation in T0 Theory}
	
	\subsection{Fundamental Mass Quantization Principle}
	
	In T0 theory, particle masses are \textbf{quantized} and follow from the fundamental geometry parameter $\xi$ through discrete scaling relationships:
	
	\begin{equation}
		m_i^{\mathrm{T0}} = n_i \cdot Q_m^{\mathrm{T0}} \cdot f_i(\xi)
	\end{equation}
	
	where:
	\begin{itemize}
		\item $n_i \in \mathbb{N}$ - Quantum number (discrete)
		\item $Q_m^{\mathrm{T0}}$ - Fundamental mass quantum in T0 units
		\item $f_i(\xi)$ - Particle-specific geometry function
	\end{itemize}
	
	\subsection{Electron Mass as Reference}
	
	The electron mass serves as the fundamental reference mass:
	
	\begin{align}
		\xi_e &= \frac{4}{3} \times 10^{-4} \times f_e(1,0,1/2) \\
		m_e^{\mathrm{T0}} &= Q_m^{\mathrm{T0}} \cdot \frac{\xi}{\xi_e} = 0.511
	\end{align}
	
	\subsection{Complete Particle Mass Spectrum}
	
	For detailed derivations of all elementary particle masses within the T0 framework, including quarks, leptons, and gauge bosons, refer to the separate comprehensive treatment ``Particle Masses in T0 Theory'' which provides:
	
	\begin{itemize}
		\item Complete mass calculations for all Standard Model particles
		\item Derivation of mass quantization rules
		\item Explanation of generation patterns
		\item Comparison with experimental values
		\item Fractal correction procedures for precision matching
	\end{itemize}
	
	\section{Important: Explicit SI Units are Necessary for\dots}
	\label{sec:si-notwendig}
	
	\subsection{1. Experimental Verification}
	Every measurement is performed in SI units:
	\begin{itemize}
		\item Particle masses in MeV/c²
		\item Cross sections in barn
		\item Magnetic moments in $\mu_B$
	\end{itemize}
	
	\subsection{2. Technological Applications}
	\begin{itemize}
		\item Detector design (lengths in m, times in s)
		\item Accelerator technology (energies in eV)
		\item Medical physics (dosage measurements)
	\end{itemize}
	
	\subsection{3. Interdisciplinary Communication}
	\begin{itemize}
		\item Astrophysics (redshifts, Hubble constant)
		\item Materials science (lattice constants)
		\item Engineering
	\end{itemize}
	
	\section{Concrete Conversion in T0 Theory}
	\label{sec:umrechnung}
	
	\subsection{Example: Electron Mass}
	\textbf{In T0 geometric units:}
	\begin{equation}
		m_e^{\mathrm{T0}} = 0.511 \quad \text{(electron mass in MeV)}
	\end{equation}
	\textbf{In SI units:}
	\begin{equation}
		m_e^{\mathrm{SI}} = m_e^{\mathrm{T0}} \cdot S_{T0} = 0.511 \cdot 1.782662 \times 10^{-30} = 9.1093837 \times 10^{-31}~\mathrm{kg}
	\end{equation}
	
	\subsection{The Fundamental Scaling Relationship}
	The conversion from T0 geometric quantities to SI units is accomplished by:
	\begin{equation}
		[\mathrm{SI}] = [\mathrm{T0}] \times S_{\text{T0}}
	\end{equation}
	where $S_{\text{T0}} = 1.782662 \times 10^{-30}$ is the scaling factor from Section~\ref{sec:scaling-derivation}; by construction it equals $1~\mathrm{MeV}/c^2$.
	
	\section{Correct Energy Scale for the Fine Structure Constant}
	
	The fundamental relationship for the fine structure constant requires a precise energy reference:
	
	\begin{align}
		\alpha &= \xi \cdot \left( \frac{E_0}{1~\mathrm{MeV}} \right)^2 \\
		\text{with} \quad E_0 &= 7.398~\mathrm{MeV} = 7.348/\sqrt{K_{\text{frak}}} \quad \text{(characteristic energy, fractally corrected)}
	\end{align}
	
	This yields:
	\begin{align}
		\alpha &= 1.333333 \times 10^{-4} \cdot (7.398)^2 \\
		&= 1.333333 \times 10^{-4} \cdot 54.73 \\
		&= 7.297 \times 10^{-3} \\
		\frac{1}{\alpha} &= 137.035
	\end{align}
	
	in agreement with the experimental value $1/\alpha = 137.036$. The rounded value $E_0 = 7.400$~MeV gives $7.3013 \times 10^{-3}$, i.e.\ $1/\alpha = 136.96$; without fractal correction ($E_0 = \sqrt{m_e m_\mu} = 7.348$~MeV) one obtains $1/\alpha = 138.9$.
	
	\section{Integration of Fractal Correction into Natural Units}
	
	The formulas in T0 theory fit in natural units without explicit fractal correction, because these units isolate the geometric essence of the theory. For exact conversions to SI units, however, fractal correction is essential to incorporate self-similar corrections of the vacuum geometry.
	
	\subsection{Why Do the Formulas Fit in Natural Units Without Fractal Correction?}
	
	In natural units, physics is reduced to a geometric, dimensionless basis (cf. Section~\ref{sec:grundprinzip}). The fundamental constants serve only as a scale, and the core formulas hold approximately without additional corrections because:
	
	\begin{itemize}
		\item \textbf{The $\xi$-parameter is intrinsically dimensionless}: $\xi$ represents the pure geometry of the vacuum field and acts like a ``universal scaling factor.''
		
		\item \textbf{Approximate validity for rough calculations}: Many T0 formulas are exact in the geometric ideal form, without correction.
		
		\item \textbf{Example: Electron mass in natural units}:
		\begin{equation}
			m_e^{\mathrm{T0}} = 0.511 \quad \text{(electron mass in MeV)}
		\end{equation}
		This ``fits'' immediately because $0.511$ is the measured value in MeV; the corpus ladder $m_e = \frac43\,\xi^{3/2}\,v$ gives $0.505$ with $v$ in MeV without correction.
	\end{itemize}
	
	\subsection{Why is Fractal Correction Necessary for Exact SI Conversions?}
	
	SI units are human conventions that ``contaminate'' the geometric purity of T0 theory. To achieve exact agreement with experiments, fractal correction must be \textbf{explicitly applied} because:
	
	\begin{itemize}
		\item \textbf{Fractal self-similarity breaks scale invariance}
		\item \textbf{Conversion requires explicit scaling}
		\item \textbf{Cosmological reference effects}
	\end{itemize}
	
	\subsection{Mathematical Specification of Fractal Correction}
	
	The fractal correction is explicitly defined as:
	\begin{equation}
		f_{\text{fractal}}(E_0) = \prod_{n=1}^{137} \left(1 + \delta_n \cdot \xi \cdot \left(\frac{4}{3}\right)^{n-1}\right)
	\end{equation}
	where $\delta_n$ are dimensionless coefficients describing the fractal structure at each stage.
	
	\subsection{Comparison: Approximation vs. Exactness}
	
	\begin{table}[h]
		\centering
				\adjustbox{max width=\textwidth}{%
\begin{tabular}{p{4cm}p{6cm}p{6cm}}
			\toprule
			\textbf{Aspect} & \textbf{Without fractal correction (T0 units)} & \textbf{With fractal correction (for SI conversion)} \\
			\midrule
			Accuracy & Approximate ($\sim 98$--$99$\,\%, geometrically ideal) & Exact (to $10^{-6}$, matches CODATA measurements) \\
			Example: $\alpha$ & $\alpha \approx \xi \cdot (E_0)^2 \approx 1/137$ (rough) & $\alpha = 1/137.03599\dots$ (via 137 stages) \\
			Mass calculation & $m_e^{\mathrm{T0}} = 0.511$ (in MeV) & $m_e^{\mathrm{SI}} = 9.1093837\times 10^{-31}$ kg (physical) \\
			Energy scale & $E_0 = \sqrt{m_e m_\mu} = 7.348$ MeV (ideal, bare mean without $K_{\text{frak}}$) & $E_0 = 7.398$ MeV ($7.348/\sqrt{K_{\text{frak}}}$) \\
			Scaling factor & $S_{T0} = 1.782662\times 10^{-30}$ (by construction) & $S_{T0} \cdot R_f$ (renormalized) \\
			Advantage & Fast, transparent calculations & Testability with experiments \\
			Disadvantage & Ignores fractal subtleties & Complex (iteration over resonance stages) \\
			\bottomrule
		\end{tabular}%
}
		\caption{Comparison of geometric idealization in T0 units and physical exactness with fractal correction.}
		\label{tab:approximation-exaktheit}
	\end{table}
	
	\section{Important Conceptual Clarifications}
	
	When applying T0 theory, note these fundamental distinctions:
	
	\begin{itemize}
		\item \textbf{T0 quantities} are given in T0 units (e.g., $m_e^{\mathrm{T0}} = 0.511$, the electron mass in MeV)
		\item \textbf{SI quantities} are physical measurements (e.g., $m_e^{\mathrm{SI}} = 9.1093837\times 10^{-31}$ kg)
		\item \textbf{$S_{T0}$} is the scaling between these realms; $S_{T0} = 1~\mathrm{MeV}/c^2$ holds by construction
		\item The energy reference for $\alpha$ is $E_0 = 7.398$ MeV $= 7.348/\sqrt{K_{\text{frak}}}$; in the geometric idealization without $K_{\text{frak}}$ it is $\sqrt{m_e m_\mu} = 7.348$ MeV
		\item All mass scales are \textbf{discretely quantized} in both T0 and SI representations
	\end{itemize}
	
	\section{Special Significance for T0 Theory}
	
	\subsection{The Deeper Insight}
	T0 theory reveals that natural units are not merely a calculational convenience, but express the \textbf{true geometric nature of physics}:
	\begin{itemize}
		\item \textbf{$\xi$} is the fundamental dimensionless geometry constant
		\item \textbf{$S_{T0}$} connects geometric idealization to physical measurement
		\item \textbf{T0 quantities} represent the ideal geometric forms
		\item \textbf{SI quantities} are their measurable projections into our physical reality
		\item \textbf{Particle masses} are quantized geometric patterns in both realms
	\end{itemize}
	
	\subsection{Practical Implications}
	\begin{enumerate}
		\item \textbf{Theoretical development}: Work in T0 units using geometric quantities
		\item \textbf{Fundamental scaling}: Apply $S_{T0}$ to project to physical reality
		\item \textbf{Predictions}: Convert to SI units for experimental verification
		\item \textbf{Verification}: Compare with measured SI values
		\item \textbf{Quantization}: Respect the discrete nature of all physical scales
	\end{enumerate}
	
	\section{Notation and Symbols}
	
	\begin{table}[h]
		\centering
		\adjustbox{max width=\textwidth}{%
\begin{tabular}{p{3cm}p{10cm}}
			\toprule
			\textbf{Symbol} & \textbf{Meaning and Explanation} \\
			\midrule
			$c$ & Speed of light in vacuum; fundamental constant of nature \\
			$\hbar$ & Reduced Planck constant \\
			$k_B$ & Boltzmann constant \\
			$G$ & Gravitational constant \\
			$E$ & Energy; in natural units dimensionally equivalent to mass and frequency \\
			$m$ & Mass; in natural units $m = E$ (since $c=1$) \\
			$p$ & Momentum; in natural units dimensionally equivalent to energy \\
			$\omega$ & Angular frequency; in natural units $\omega = E$ (since $\hbar=1$) \\
			$\alpha$ & Fine structure constant; dimensionless coupling constant \\
			$\xi$ & Fundamental geometry parameter of T0 theory; $\xi = \frac{4}{3} \times 10^{-4}$ \\
			$E_0$ & Reference energy in T0 theory; $E_0 = 7.398~\mathrm{MeV} = 7.348/\sqrt{K_{\text{frak}}}$ \\
			$m_e^{\mathrm{T0}}$ & Electron mass in T0 units; $m_e^{\mathrm{T0}} = 0.511$ (electron mass in MeV) \\
			$m_e^{\mathrm{SI}}$ & Electron mass in SI units; $m_e^{\mathrm{SI}} = 9.1093837\times 10^{-31}$ kg (physical) \\
			$[E]$ & Energy dimension; fundamental dimension in natural units \\
			SI & International System of Units (physical measurements) \\
			T0 & T0 geometric units (ideal geometric forms) \\
			$S_{T0}$ & Scaling factor; $S_{T0} = 1.782662 \times 10^{-30}$ (by construction $1~\mathrm{MeV}/c^2$ in kg) \\
			$R_f$ & Fractal correction factor \\
			$f_{\text{fractal}}$ & Fractal correction function \\
			$Q_m^{\mathrm{T0}}$ & Fundamental mass quantum in T0 units \\
			$Q_m^{\mathrm{SI}}$ & Fundamental mass quantum in SI units \\
			$n_i$ & Quantum number for particle $i$; $n_i \in \mathbb{N}$ (discrete) \\
			$\delta_n$ & Fractal correction coefficients; dimensionless \\
			\bottomrule
		\end{tabular}%
}
		\caption{Explanation of the notation and symbols used}
	\end{table}
	
	\section{Fundamental Relationships}
	
	\begin{table}[h]
		\centering
		\adjustbox{max width=\textwidth}{%
\begin{tabular}{p{4cm}p{10cm}}
			\toprule
			\textbf{Relationship} & \textbf{Meaning} \\
			\midrule
			$E = m$ & Mass-energy equivalence (since $c=1$) \\
			$E = \omega$ & Energy-frequency relationship (since $\hbar=1$) \\
			$[L] = [T] = [E]^{-1}$ & Length and time have same dimension as inverse energy \\
			$[m] = [p] = [E]$ & Mass and momentum have same dimension as energy \\
			$\alpha = \xi (E_0/1\mathrm{MeV})^2$ & Fundamental relationship in T0 theory \\
			$m_i^{\mathrm{T0}} = n_i \cdot Q_m^{\mathrm{T0}} \cdot f_i(\xi)$ & Quantized mass formula in T0 units \\
			$m_i^{\mathrm{SI}} = m_i^{\mathrm{T0}} \cdot S_{T0}$ & Fundamental scaling to SI units \\
			$S_{T0} = \dfrac{m_e^{\mathrm{SI}}}{m_e^{\mathrm{T0}}}$ & Definition of fundamental scaling factor \\
			\bottomrule
		\end{tabular}%
}
		\caption{Fundamental relationships in T0 theory and scaling to physical units}
	\end{table}
	
	\section{Conversion Factors}
	
	\begin{table}[h]
		\centering
		\adjustbox{max width=\textwidth}{%
\begin{tabular}{lll}
			\toprule
			\textbf{Quantity} & \textbf{Conversion Factor} & \textbf{Value} \\
			\midrule
			$S_{T0}$ & Scaling factor & $1.782662 \times 10^{-30}$ \\
			$m_e^{\mathrm{T0}}$ & Electron mass (T0 units) & $0.511$ \\
			$m_e^{\mathrm{SI}}$ & Electron mass (SI units) & $9.1093837 \times 10^{-31}~\mathrm{kg}$ \\
			$1~\mathrm{MeV}/c^2$ & Conventional mass unit & $1.782662 \times 10^{-30}~\mathrm{kg}$ \\
			$1~\mathrm{MeV}$ & Energy in joules & $1.602176 \times 10^{-13}~\mathrm{J}$ \\
			$1~\mathrm{fm}$ & Length in natural units & $5.06773 \times 10^{-3}~\mathrm{MeV}^{-1}$ \\
			\bottomrule
		\end{tabular}%
}
		\caption{Fundamental conversion factors between T0 geometric units and SI physical units}
	\end{table}
